\section*{第 \ref{ch:b3:seismology-imaging} 章——地震学与地球成像}

\begin{solution}{exo:b3:seismology-imaging:1}
(a) $\sin i_0 = pv_0/R = 500\times5.80/6371 = 0.4552$，故 $i_0 = \ang{27.1}$。
(b) 在 \qty{771}{km} 处，$r = \qty{5600}{km}$，$r/v = 5600/11.066 =
\qty{506.1}{s.rad^{-1}}$；在 \qty{871}{km} 处，$r = \qty{5500}{km}$，$r/v =
5500/11.245 = \qty{489.1}{s.rad^{-1}}$。射线在 $r/v = p = 500$ 处转折：
$z_t = 771 + 100\times(506.1 - 500)/(506.1 - 489.1) = 771 + 36 =
\qty{807}{km}$。
\end{solution}

\begin{solution}{exo:b3:seismology-imaging:2}
(a) $(455.5 - 369.6)/10 = \qty{8.59}{s}$ 每度，且
$8.59\times180/\pi = \qty{492.2}{s.rad^{-1}}$。(b) 一度弧长为
$6371\pi/180 = \qty{111.19}{km}$，故波前扫过地表的速度为
$111.19/8.59 = \qty{12.94}{km.s^{-1}}$，远快于地表任何岩石中的波速，
射线几乎竖直到达。(c) 在转折点处 $v = r_t/p =
5450/492.2 = \qty{11.07}{km.s^{-1}}$，为下地幔 P 波速度。
\end{solution}

\begin{solution}{exo:b3:seismology-imaging:3}
(a) $\mu_1 = \rho_1\beta_1^2 = 2700\times3500^2 = \qty{3.31e10}{Pa}$
（\qty{33.1}{GPa}）；$\mu_2 = 3300\times4500^2 = \qty{6.68e10}{Pa}$
（\qty{66.8}{GPa}）。(b) $\lambda = cT = 3.79\times20 = \qty{75.8}{km}$，为层厚的 2.17
倍：波深入半空间很深，这正是
其速度已高于 $\beta_1$ 的原因。(c) $4000/3.384 = \qty{1182}{s}$，且
$4000/4.333 = \qty{923}{s}$：20 秒波比 100 秒波晚到 \qty{259}{s}，
即 \qty{4.3}{min}。
\end{solution}

\begin{solution}{exo:b3:seismology-imaging:4}
(a) ${}_0S_2$：$1/(53.89\times60) = \qty{3.093e-4}{Hz} = \qty{0.3093}{mHz}$；
${}_0T_2$：$1/(44.01\times60) = \qty{0.3787}{mHz}$。(b) 每天的振荡次数分别为 $1440/53.89 = 26.7$
与 $1440/44.01 = 32.7$。(c) 两者之差为
$\qty{0.0694}{mHz}$，其倒数为 \qty{14400}{s}：至少需要 \qty{4.0}{h} 的
记录，实际上需要更长，因为衰减使谱峰展宽并
使其分裂。
\end{solution}

\begin{solution}{exo:b3:seismology-imaging:5}
(a) 球面上两点之间、对球心张角为 $\Delta$ 的弦，
其长度为 $2R\sin(\Delta/2)$，以速度 $v$ 传播。
(b) 在 \ang{30} 处：$2\times6371\times\sin\ang{15}/10.0 = \qty{329.8}{s}$，对比
\qty{369.6}{s}；在 \ang{90} 处：$2\times6371\times\sin\ang{45}/10.0 =
\qty{901.0}{s}$，对比 \qty{779.7}{s}。(c) 到 \ang{30} 的射线在
上地幔中转折，那里的真实速度约为 8 至 \qty{10}{km.s^{-1}}，低于
均匀模型的 \qty{10}{km.s^{-1}}；到 \ang{90} 的射线大部分路径位于
下地幔中，速度为 11 至 \qty{13.7}{km.s^{-1}}，并向下弯入其中，故
真实地球更快。
\end{solution}

\begin{solution}{exo:b3:seismology-imaging:6}
(a) 两边平方，$(2 - \xi)^4 = 16(1 - \xi/3)(1 - \xi)$。展开：
$16 - 32\xi + 24\xi^2 - 8\xi^3 + \xi^4 = 16 - \tfrac{64}{3}\xi + \tfrac{16}{3}\xi^2$，
故 $\xi^4 - 8\xi^3 + \tfrac{56}{3}\xi^2 - \tfrac{32}{3}\xi = 0$，即
$\xi\,(3\xi^3 - 24\xi^2 + 56\xi - 32) = 0$，且
$3\xi^3 - 24\xi^2 + 56\xi - 32 = (\xi - 4)(3\xi^2 - 12\xi + 8)$。根
$\xi = 0$ 不是波。(b) 二次式的根为 $2 \pm 2/\sqrt3$，
即 3.155 与 0.8453；只有 $\xi = 0.8453$ 位于 $(0, 1)$ 内（$\xi = 4$
不在其中），且它满足未平方的方程，两边均为 $4/3$。
故 $c/\beta = \sqrt{0.8453} = 0.9194$。(c) $c = 0.9194\times3.46 =
\qty{3.181}{km.s^{-1}}$；在 \qty{1000}{km} 处，$314.4 - 289.0 = \qty{25.3}{s}$
晚于 S 波到达。
\end{solution}

\begin{solution}{exo:b3:seismology-imaging:7}
$\mathsf G$ 的各行为 $(1, x_i)$，$\mathsf C_d^{-1} = \operatorname{diag}
(100, 100, 25, 25)$（\unit{s^{-2}}）。正规方程为
\[
  \begin{pmatrix} 250 & 23750 \\ 23750 & \num{2812500} \end{pmatrix}
  \begin{pmatrix} a \\ b \end{pmatrix}
  = \begin{pmatrix} 3665 \\ 418125 \end{pmatrix},
\]
行列式为 \num{1.390625e8}。于是 $b = (250\times418125 -
23750\times3665)/\num{1.390625e8} = \qty{0.12575}{s.km^{-1}}$，且
$a = (2812500\times3665 - 23750\times418125)/\num{1.390625e8} =
\qty{2.713}{s}$。协方差为 $\mathsf N^{-1}$：$\sigma_b =
\sqrt{250/\num{1.390625e8}} = \qty{0.00134}{s.km^{-1}}$，$\sigma_a =
\sqrt{2812500/\num{1.390625e8}} = \qty{0.142}{s}$。速度为
$1/b = \qty{7.95}{km.s^{-1}}$，误差为 $\sigma_b/b^2 =
\qty{0.08}{km.s^{-1}}$：这是上地幔最顶部的速度。残差为 $-0.001$、
$0.011$、$-0.076$ 与 $\qty{0.036}{s}$，且 $\chi^2 = 0.19$，对应 $4 - 2 = 2$
个自由度：小于预期，故所给误差若有偏差，
则偏于保守；拟合是可以接受的。
\end{solution}

\begin{solution}{exo:b3:seismology-imaging:8}
(a) $\mathsf G = (30\;\;40)$ 为单行矩阵：$\mathsf G^{\mathsf T}\mathsf G$ 为
秩为一的 $2\times2$ 矩阵，故奇异；任何满足 $30m_1 + 40m_2 = 0$ 的模型都可
加入而不改变数据。(b) $\mathsf G\mathsf G^{\mathsf T} =
900 + 1600 = 2500$，故 $\mathbf m = (30, 40)\times0.50/2500 =
(0.006, 0.008)\,\unit{s.km^{-1}}$。(c) $\mathsf R = \frac{1}{2500}
\begin{pmatrix} 900 & 1200 \\ 1200 & 1600 \end{pmatrix} =
\begin{pmatrix} 0.36 & 0.48 \\ 0.48 & 0.64 \end{pmatrix}$，迹为 1：一个
数据只能分辨一种参数组合。(d) $\mathsf R\,(0.01, 0) =
(0.0036, 0.0048)\,\unit{s.km^{-1}}$：异常沿射线分布，在射线较长的
单元中分布更多。
\end{solution}

\begin{solution}{exo:b3:seismology-imaging:9}
(a) 线性化：$\delta t = -(\delta v/v)L/v = -0.03\times200/9.0 =
\qty{-0.667}{s}$；精确值为 $L/(1.03v) - L/v = 200/9.27 - 200/9.0 =
\qty{-0.647}{s}$。(b) 射线在板片中的路径为 $100/\cos\ang{40} =
\qty{130.5}{km}$：线性化时 $-0.03\times130.5/9.0 = \qty{-0.435}{s}$，
精确值为 $\qty{-0.422}{s}$。(c) 精确值等于线性化结果除以
$1.03$：误差为 \qty{3}{\percent}，与 $\delta v/v$ 本身同量级，符合
一阶理论的预期。
\end{solution}

\begin{solution}{exo:b3:seismology-imaging:10}
(a) 该泛频在其半空间部分衰减时存在，即 $c < \beta_2$；在其
截止处 $c = \beta_2$，故 $q_2 = 0$，且 $\tan(kHq_1) = 0$，其中 $kHq_1$ 位于
$[n\pi, n\pi + \pi/2)$ 内：$kHq_1 = n\pi$。(b) 取 $k = \omega/\beta_2$ 及
$q_1 = \sqrt{\beta_2^2/\beta_1^2 - 1} = \sqrt{\beta_2^2 - \beta_1^2}/\beta_1$，
$\omega_nH\sqrt{\beta_2^2 - \beta_1^2}/(\beta_1\beta_2) = n\pi$，即
所求公式。(c) $\sqrt{4.5^2 - 3.5^2} = \sqrt8 = 2.828$；$\omega_1 =
\pi\times3.5\times4.5/(35\times2.828) = \qty{0.4998}{rad.s^{-1}}$，周期为
$2\pi/0.4998 = \qty{12.57}{s}$：在更长周期上只存在
基模。(d) 当 $n = 0$ 时截止频率为零：基模
在所有周期上都存在。短周期下 $kH \to \infty$，而 $kHq_1$ 保持
小于 $\pi/2$，故 $q_1 \to 0$ 且 $c \to \beta_1$。
\end{solution}

\begin{solution}{exo:b3:seismology-imaging:11}
(a) $\mathsf G^{\mathsf T}(\mathsf G\mathsf G^{\mathsf T} + \varepsilon^2\mathsf I)
= (\mathsf G^{\mathsf T}\mathsf G + \varepsilon^2\mathsf I)\mathsf G^{\mathsf T}$;
在左侧乘以 $(\mathsf G^{\mathsf T}\mathsf G + \varepsilon^2
\mathsf I)^{-1}$，右侧乘以 $(\mathsf G\mathsf G^{\mathsf T} +
\varepsilon^2\mathsf I)^{-1}$，两者在 $\varepsilon > 0$ 时均可逆。
(b) 此处 $\mathsf G\mathsf G^{\mathsf T} = \|\mathbf g\|^2 = 2500$ 是一个数，
故 $\mathsf R_\varepsilon = \mathsf G^{\mathsf T}\mathsf G/(2500 + \varepsilon^2)$，
其迹为 $2500/(2500 + \varepsilon^2)$。(c) 预测数据为
$\mathbf g\cdot\mathbf m_\varepsilon = 0.50\times2500/(2500 + \varepsilon^2)$
残差为 $0.50\,\varepsilon^2/(2500 + \varepsilon^2)$。
$\varepsilon^2 = 0$ 时：迹为 1，残差为 0；对于 2500：迹为 0.5，残差为
\qty{0.25}{s}；对于 \num{10000}：迹为 0.2，残差为 \qty{0.40}{s}。阻尼越大，
模型更小，分辨率更差，拟合差更大。
\end{solution}

\begin{solution}{exo:b3:seismology-imaging:12}
(a) 由练习 5 可得 $T = (2R/v)\sin(\Delta/2)$，故 $\mathrm dT/\mathrm d\Delta
= (R/v)\cos(\Delta/2)$。圆心处张角为 $\Delta$ 的弦，
是顶角为 $\Delta$ 的等腰三角形的底边，故它在震源处与半径所成的角为
$i_0 = (\pi - \Delta)/2$，且 $p = R\sin i_0/v =
(R/v)\cos(\Delta/2)$，二者相同。(b) 转折点位于
弦的中点，距圆心 $R\cos(\Delta_1/2)$；等价地，
$r_1 = p_1v$。(c) 此处 $p(\Delta)/p_1 = \cos(\Delta/2)/\cos(\Delta_1/2)$，且
由公式得 $r_1 = R\exp\bigl(-\ln(1/\cos(\Delta_1/2))\bigr) =
R\cos(\Delta_1/2)$，与 (b) 一致；速度为 $r_1/p_1 = R\cos(\Delta_1/2)\big/
\bigl((R/v)\cos(\Delta_1/2)\bigr) = v$.
\end{solution}

\begin{solution}{pb:b3:seismology-imaging:1}
\textbf{1.} $i_0 = \arcsin 0.8 = \ang{53.13}$，且 $p = 0.8/5.0 =
\qty{0.16}{s.km^{-1}}$。

\textbf{2.} 射线在 $v = 1/p = \qty{6.25}{km.s^{-1}}$ 处转折，即在
$z = (6.25 - 5.0)/0.05 = \qty{25}{km}$ 处。

\textbf{3.} 沿射线对 $\sin i = pv$ 求导，得 $\cos i\,\mathrm di/
\mathrm ds = pg\,\mathrm dz/\mathrm ds = pg\cos i$，故 $\mathrm di/\mathrm ds =
pg$：射线方向以恒定速率转动，恰如半径为
$1/(pg) = 1/(0.16\times0.05) = \qty{125}{km}$ 的圆的曲率。在转折点处射线
水平，位于 \qty{25}{km} 深度，圆心竖直位于其上方 \qty{125}{km}
：在 $z = \qty{-100}{km}$ 处，即地表以上 $v_0/g$ 处，此处
外推速度将为零。

\textbf{4.} 在震源处，圆的半径与射线垂直，
故它与水平面的夹角为 $i_0$，圆心水平距离震源
$125\cos i_0$，由对称性
$X = 2\times125\times\cos i_0 = 2\times125\times0.6 = \qty{150}{km}$，或
$X = 2\sqrt{1 - p^2v_0^2}/(pg)$。

\textbf{5.} $\mathrm dT = \mathrm ds/v$，其中 $\mathrm ds = \mathrm di/(pg)$，且
$v = \sin i/p$：$\mathrm dT = \mathrm di/(g\sin i)$。于是
\[
  T = \frac2g\int_{i_0}^{\pi/2}\frac{\mathrm di}{\sin i} = \frac2g\ln\cot\frac{i_0}2 =
  \frac2g\operatorname{arccosh}\frac1{pv_0} .
\]
此处 $\cot(\ang{26.57}) = 2$，故
$T = 40\ln2 = \qty{27.73}{s}$。

\textbf{6.} 当 $p = 0.15$ 时：$pv_0 = 0.75$，$\sqrt{1 - 0.5625} = 0.6614$，
$X = 2\times0.6614/0.0075 = \qty{176.4}{km}$，$T = 40\operatorname{arccosh}
(1.3333) = 40\times0.7954 = \qty{31.81}{s}$。于是 $\Delta T/\Delta X =
4.08/26.4 = \qty{0.155}{s.km^{-1}}$，介于 0.15 与 0.16 之间，与
$\mathrm dT/\mathrm dX = p$ 的要求相符。

\textbf{7.} $\mu_1 = 2700\times3500^2 = \qty{33.1}{GPa}$，且
$\mu_2 = 3300\times4500^2 = \qty{66.8}{GPa}$。

\textbf{8.} $k = 2\pi/(cT) = 2\pi/(3.99\times25) = \qty{0.06299}{rad.km^{-1}}$；
$q_1 = \sqrt{(3.99/3.5)^2 - 1} = 0.5474$；$q_2 = \sqrt{1 - (3.99/4.5)^2} = 0.4624$。

\textbf{9.} $\tan(kHq_1) = 66.825\times0.4624/(33.075\times0.5474) = 1.707$；基模
取 $kHq_1 = \arctan1.707 = 1.041$，故
$H = 1.041/(0.06299\times0.5474) = \qty{30.2}{km}$。

\textbf{10.} 在 \qty{40}{s} 时：$k = 2\pi/170 = 0.03696$，$q_1 = 0.6888$，
$q_2 = 0.3287$，$\tan(kHq_1) = 0.964$，$kHq_1 = 0.7671$，且
$H = 0.7671/(0.03696\times0.6888) = \qty{30.1}{km}$。两者在
相速度的精度范围内一致：地壳厚约 \qty{30}{km}。

\textbf{11.} $\omega = 0.2513$ 与 $\qty{0.1571}{rad.s^{-1}}$，$k = 0.06299$ 与
$\qty{0.03696}{rad.km^{-1}}$：$U \approx 0.0942/0.02603 = \qty{3.62}{km.s^{-1}}$，
小于两个相速度，这在 $c$ 随
周期增大时必然成立。

\textbf{12.} $\lambda = 3.99\times25 = \qty{99.8}{km}$ 与 $4.25\times40 =
\qty{170}{km}$。基模的采样深度约为其波长的一部分：
25 秒的波主要位于地壳中，40 秒的波则深入地幔较多，
因此其速度更接近 $\beta_2$。两个周期确定相同的
$H$，说明双层模型是合适的。

\textbf{13.} 以单元 1 至 4 为列、A 至 F 为行，单位为
\qty{1000}{km}（$\sqrt2\times100/1000 = 0.1414$）：
\[
  \mathsf G = \begin{pmatrix}
  0.1 & 0.1 & 0 & 0 \\ 0 & 0 & 0.1 & 0.1 \\ 0.1 & 0 & 0.1 & 0 \\
  0 & 0.1 & 0 & 0.1 \\ 0.1414 & 0 & 0 & 0.1414 \\ 0 & 0.1414 & 0.1414 & 0
  \end{pmatrix}.
\]

\textbf{14.} $\mathsf C_d^{-1} = \operatorname{diag}(400, 400, 400, 400, 100,
100)$。行或列射线贡献 $0.1^2\times400 = 4$，对角射线
$0.1414^2\times100 = 2$。单元 1 被 A、C（各 4）与 E（2）穿过：10。
单元 1 与 2 共享射线 A：4；单元 1 与 3 共享 C：4；单元 1 与 4 共享 E：2；
依此类推，由对称性得到所给的四行。

\textbf{15.} 行或列射线对其每个单元贡献 $0.1\times400\,\delta t = 40\,
\delta t$，对角射线贡献 $0.1414\times100\,\delta t =
14.14\,\delta t$。单元 1：$40(0.16 - 0.36) + 14.14\times(-0.15) = -10.12$；
单元 2：$40(0.16 + 0.11) + 14.14\times(-0.27) = 6.98$；单元 3：$40(-0.48 - 0.36) +
14.14\times(-0.27) = -37.42$；单元 4：$40(-0.48 + 0.11) + 14.14\times(-0.15) =
-16.92$。

\textbf{16.} 对 $(1,1,1,1)$，$\mathsf N$ 的每一行之和均为 20。对
$(1,-1,-1,1)$，第 1 行给出 $10 - 4 - 4 + 2 = 4$，其余各行
模式相同。对 $(1,1,-1,-1)$，$10 + 4 - 4 - 2 = 8$，对 $(1,-1,1,-1)$，
$10 - 4 + 4 - 2 = 8$。这些向量两两正交且范数为 2，因此
$\mathsf N^{-1} = \sum_k \mathbf u_k\mathbf u_k^{\mathsf T}/(4\lambda_k)$；其
$(1,1)$ 元素为 $\tfrac14(\tfrac1{20} + \tfrac14 + \tfrac18 + \tfrac18) =
\tfrac{11}{80}$，$(1,2)$ 元素为 $\tfrac14(\tfrac1{20} - \tfrac14 + \tfrac18
- \tfrac18) = -\tfrac4{80}$，$(1,4)$ 元素为 $\tfrac14(\tfrac1{20} +
\tfrac14 - \tfrac18 - \tfrac18) = \tfrac1{80}$，依此类推。

\textbf{17.} $\hat{\mathbf m} = \mathsf N^{-1}\mathbf b$：
$\hat m_1 = (11\times(-10.121) - 4\times6.982 - 4\times(-37.418) - 16.921)/80
= -0.081$；同理 $\hat m_2 = 1.844$，$\hat m_3 = -3.706$，
以及 $\hat m_4 = \qty{-0.931}{ms.km^{-1}}$。

\textbf{18.} $\delta v/v_0 = -v_0\,\delta s = -8.0\times10^{-3}\,\hat m$，以百分数表示
$-0.8\,\hat m$：单元 1 为 $+0.07$，单元 2 为 $-1.48$，单元 3 为 $+2.96$，单元 4
$\qty{+0.75}{\percent}$。

\textbf{19.} $\mathsf N^{-1}$ 的每个对角元素均为 $11/80$：每个单元的
标准误差为 $\sqrt{0.1375} = \qty{0.371}{ms.km^{-1}}$，即
$\qty{0.30}{\percent}$。单元 1 和单元 2 的相关系数为 $-4/11 = -0.36$：
两者共享射线 A，且在一个单元中引入的慢度可与另一单元相互抵换。
单元 1 与单元 4 仅共享精度较低的对角射线 E，二者
呈弱正相关，$+1/11 = +0.09$。

\textbf{20.} 预测残差分别为 $0.176$、$-0.464$、$-0.379$、$0.091$，
$-0.143$ 与 $\qty{-0.263}{s}$；拟合后的残差为 $-0.016$，
$-0.016$、$0.019$、$0.019$、$-0.007$ 与 $\qty{-0.007}{s}$，
$\chi^2 = 0.50$，自由度为 $6 - 4 = 2$：与所给
误差相符，若有偏差则略偏小。

\textbf{21.} 特征向量与 $\mathsf N$ 的相同，本征值为
$\lambda/(\lambda + 4)$：$20/24 = 5/6$，$4/8 = 1/2$，$8/12 = 2/3$（出现两次）。与
第 16 题一样，$R_{11} = \tfrac14(\tfrac56 + \tfrac12 + \tfrac23 + \tfrac23) =
\tfrac23$；$R_{12} = \tfrac14(\tfrac56 - \tfrac12 + \tfrac23 - \tfrac23) =
\tfrac1{12}$；$R_{13} = \tfrac14(\tfrac56 - \tfrac12 - \tfrac23 + \tfrac23) =
\tfrac1{12}$；$R_{14} = \tfrac14(\tfrac56 + \tfrac12 - \tfrac23 - \tfrac23) = 0$。

\textbf{22.} 将 $\mathbf b$ 投影到各特征向量上：$\mathbf b\cdot\mathbf u_k/4
= -14.37$、$0.849$、$12.80$ 与 $-9.40$，分别对应本征值 20、4、8 和 8。
阻尼估计将它们除以 $\lambda + 4$ 而非 $\lambda$；对
单元 3（分量为 $1, -1, -1, 1$），
$m_3 = -14.37/24 - 0.849/8 - 12.80/12 - 9.40/12 = \qty{-2.555}{ms.km^{-1}}$，
即速度异常为 $\qty{+2.04}{\percent}$。阻尼使每个分量都
趋于零，其中约束最弱的分量受影响最大。

\textbf{23.} 未阻尼模型在误差范围内拟合数据，
（$\chi^2 = 0.50$，自由度为 2），而阻尼模型对数据的拟合
很差（$15.8$）。仅当 $\mathsf N$ 奇异或接近奇异时
才需要阻尼；此处最小本征值 4 并不小，估计稳定，且
其误差适中。

\textbf{24.} 第 3 行为 $(1/12, 0, 2/3, 1/12)$：单元 3 的阻尼估计
等于其真值的三分之二，再加上单元 1 和单元 4 的各十二分之一，它们是
沿射线 C 与射线 B 的相邻单元，而不含单元 2 的任何成分。

\textbf{25.} 单元 3（左侧下方的单元）约快
\qty{3}{\percent}；单元 2 偏慢。冷而致密的大洋岩石圈板片
下沉于大陆之下，比周围地幔快，因此
位于 K 洋侧的深部快速单元即为板片；其上方另一侧的慢速单元
它位于另一侧，正处于地幔楔较热之处。

\textbf{26.} 由第 18、19 题可得，板片单元的 P 波速度异常为
\textbf{\qty[text-series-to-math=true, uncertainty-mode=separate]{3.0 \pm 0.3}{\percent}}。
\end{solution}
